\subsection{An abstract result}\label{s:abstrIntro} Even though we have presented several results related to the semilinear wave and plate equations, at the core of all of them lies an abstract result. This result essentially states that solutions to a nonlinear problem, with a skew-adjoint linear part, and a compact nonlinearity which is also analytic, are analytic in time, provided that the observed solution is zero.

We first need to introduce some notation and assumptions in order to state the aforementioned result. Let $T>0$ and let us consider the nonlinear observability system
\begin{align}\label{intro:eq-obs}
\left\{\begin{aligned}
\partial_t U&=AU+F(U+H_1)+H_2, &t&\in (0, T), \\
\bC U(t)&=0, &t&\in (0, T),
\end{aligned}\right.
\end{align}
on a suitable real Hilbert space $X$, where $A$ is a skew-adjoint operator on $X$, $F$ is a mapping on $X$, $H_1$ and $H_2$ are some parameters, and $\bC$ is a bounded observation operator in $X$. The first assumption dictates the class of PDEs we will be working with.
\begin{assump}{1}
\label{assumAA}$A$ is a skew-adjoint operator with domain $D(A)$ on a real separable Hilbert space $X$, so that $A^*A=-A^2$ has a compact resolvent.
\end{assump}
This assumption allows us to introduce $X^\sigma$ as the interpolation space in between $D(A)$ and $X$ for $\sigma\in [0, 1]$, with $X^0=X$ and $X^1=D(A)$. From now on, we~will work at a fixed level of regularity, for which we fix $\sigma\in [0, 1]$ and instead consider \eqref{intro:eq-obs} in $X^\sigma$. We~will need to exert some control on the linear semigroup $t\in [0, T]\mto e^{At}\in \mc{L}(X^\sigma)$ generated by $A$ and even more, we~will need some extra control at a slightly higher regularity. This will be accomplished through the following assumption.

\begin{assump}{2}
\label{assumCC}Let $\bC\in\mc{L}(X^\sigma, X^\sigma)$ be an observation operator. We~assume that $t\mto e^{tA}$ is observable on $[0,T]$, namely, there exists a constant $\mathfrak{C}_{obs}>0$ such that
\begin{align}\label{observabstractintro}
\norm{W_0}_{X^\sigma}^2\leq \mathfrak{C}_{\textup{obs}}^2\int_0^T \norm{\bC e^{tA}W_0}_{X^\sigma}^2dt,\quad \forall W_0\in X^\sigma.
\end{align}
Furthermore, we~assume that $\bC\in\mc{L}(X^{\sigma+\veps}, X^{\sigma+\veps})$ for some $\veps>0$ and that there exists (another) constant $\mathfrak{C}_{obs}>0$ such that
\begin{align}\label{observabstractvepsintro}
\norm{W_0}_{X^{\sigma+\veps}}^2\leq \mathfrak{C}_{\textup{obs}}^2\int_0^T \norm{\bC e^{tA}W_0}_{X^{\sigma+\veps}}^2dt,\quad \forall W_0\in X^{\sigma+\veps}.
\end{align}
\end{assump}

We now make the last assumption regarding the nonlinearity. For a given real Banach space $Y$, we~denote the ball centered at $0$ of radius $M$ by
\begin{align*}
\mathbb{B}_M(Y):=\{y\in Y\ |\ \norm{y}_Y\leq M\}.
\end{align*}
For a given interval $I\subset \R$, we~denote the ball of radius $M$ of $C^0(I, Y)$ by
\begin{align*}
\B_{M}^I(Y)=\{U\in C^0(I, Y)\ |\ \forall t\in I,\ \norm{U(t)}_Y\leq M\},
\end{align*}
Moreover, we~introduce the canonical complexification $Y_\mathbb{C}$, defined as the set of elements $y_1+iy_2$, $y_j\in Y$, see \cite[\S 2]{BS:71-polynomials} for more details. We~then introduce the notation for the \emph{cylinder} on $Y_\mathbb{C}$
\begin{align*}
\mathbb{B}_{M, \delta}(Y)=\{y\in Y_\mathbb{C}\ |\ \norm{\Re(y)}_Y\leq M\ \text{and}\ \norm{\Im(y)}_Y\leq \delta\},
\end{align*}
and similarly on $C^0(I, Y_\mathbb{C})$
\begin{align*}
\B_{M, \delta}^I(Y)=\{U\in C^0(I, Y_\mathbb{C})\ |\ \forall t\in I,\ \norm{\Re(U(t))}_Y\leq M\ \text{and}\ \norm{\Im(U(t))}_Y\leq \delta\}.
\end{align*}
The latter space is naturally endowed with the $L^\infty(I,Y_\mathbb{C})$-norm. When working on balls not centered at $0$, we~will simply denote by $B_Y(y_0, r)$ the ball of center $y_0$ and radius $r>0$ on $Y$. With the previous notations, we~will make the following assumption on $F$.
\begin{assump}{3}
\label{assumFholom}
$F$ is a nonlinear Lipschitz and bounded operator from $\mathbb{B}_{4R_{0}}(X^{\sigma})$ into $X^{\sigma+\veps}$ for some $R_{0}>0$, $\sigma>0$ and $\veps>0$. Furthermore, there exists $\delta>0$ such that~$F$ has a holomorphic extension from $\mathbb{B}_{4R_{0},2\delta}(X^{\sigma})$ into $X^{\sigma+\veps}_{\mathbb{C}}$, namely, it~is holomorphic on the interior of $\mathbb{B}_{4R_{0},2\delta}(X^{\sigma})$ and continuous up to the boundary. Moreover, there exists $C>0$ such that
\begin{align}
\label{boundFanalyticintro}
\nor{F(U_{0})}{X^{\sigma+\veps}_{\mathbb{C}}}\leq C,\quad \nor{F(U_{0})-F(V_{0})}{X^{\sigma+\veps}_{\mathbb{C}}}\leq C \nor{U_{0}-V_{0}}{X^{\sigma}_{\mathbb{C}}}
\end{align}
for any $U_{0},~V_{0}\in \mathbb{B}_{4R_{0},2\delta}(X^{\sigma})$.
\end{assump}

Unless stated otherwise, we~will understand holomorphic extensions in the sense given in the assumption above. We~will also need a technical assumption on the pair $(A, \bC)$ related to commutator estimates, useful for the regularization of solutions.
\begin{assump}{4}
\label{assumcommu}
There exists $s>0$ such that $[(A^*A)^s,\bC]\in \mc{L}(X^{\sigma+2s}, X^{\sigma+\veps})$.
\end{assump}

We now state the main abstract result.

\begin{theorem}\label{thmabstractanalyticintro}
Let $R_{0}\!>\!0$ and $T\!>\!0$. Suppose that Assumptions \ref{assumAA}, \ref{assumCC} (with~$T$), \ref{assumFholom} (with $R_{0}$) and \ref{assumcommu} hold. Let $T^{*}\!>\!T$. Let $H_{1}\in \B_{R_{0}}^{[0,T^{*}]}(X^{\sigma})$ and $H_{2}\in C^{0}([0,T^{*}],X^{\sigma+\veps})$ that admit some extension in the space $C^{0}([0,T^{*}]+i[-\mu,\mu],X^{\sigma}_{\mathbb{C}})$, respectively in \hbox{$C^{0}([0,T^{*}]+i[-\mu,\mu],X^{\sigma+\veps}_{\mathbb{C}})$}, with $\mu\!>\!0$, so that the map
\begin{align*}
\left\{\begin{aligned}
(0,T^{*})+i(-\mu,\mu) & \to X^{\sigma}_{\mathbb{C}}\\
z&\mto H_{1}(z)
\end{aligned}\right.
\end{align*}
is holomorphic. We~assume the same for $H_{2}$ with value in $X^{\sigma+\veps}_{\mathbb{C}}$. We~assume moreover that $\Re H_{1}(z)\in \mathbb{B}_{R_0}(X^{\sigma})$ for any $z\in [0,T^{*}]+i[-\mu,\mu]$.

Then any $U\in C^{0}([0,T^{*}],X^{\sigma})$ satisfying
\begin{align}\label{UCabstractT*intro}
\left\{\begin{aligned}
\partial_t U&=AU+F(U+H_{1})+H_{2}&& \textnormal{ on } [0,T^{*}], \\
\bC U(t)&=0 &&\textnormal{ for }t\in [0,T^{*}],
\end{aligned}\right.
\end{align}
is real analytic in $t$ in $(0,T^*)$ with value in $X^{\sigma}$.
\end{theorem}
